% Source: https://artofproblemsolving.com/community/c3h2580839_doing_roots_by_hand_replicate_of_my_blog_post
% Stable identity: aops-topic-2580839-q01
\textbf{(a)} With $r=\dfrac13$:
\[\dbinom r0=1,\quad \dbinom r1=\dfrac13,\quad \dbinom r2=\dfrac{\dfrac13\cdot\left(-\dfrac23\right)}{2}=-\dfrac19,\quad \dbinom r3=\dfrac{\dfrac13\cdot\left(-\dfrac23\right)\cdot\left(-\dfrac53\right)}{6}=\dfrac5{81},\]
\[\dbinom r4=\dbinom r3\cdot\dfrac{r-3}{4}=\dfrac5{81}\cdot\left(-\dfrac23\right)=-\dfrac{10}{243}.\]

\textbf{(b)} With $x=0.8$:
\[1+\dfrac{0.8}{3}-\dfrac{0.64}{9}+\dfrac{5\cdot0.512}{81}-\dfrac{10\cdot0.4096}{243}\approx1+0.266667-0.071111+0.031605-0.016856,\]
so $\sqrt[3]{1.8}\approx\boxed{1.2103}$. (The true value is $1.21644\ldots$; more terms get closer, slowly, because $0.8$ is close to $1$.)

\textbf{(c)} For $x=9$ the ratio of consecutive terms has absolute value
\[\left|\dfrac{\dbinom r{k+1}9^{k+1}}{\dbinom rk9^k}\right|=9\cdot\dfrac{k-\dfrac13}{k+1}\ge9\cdot\dfrac13=3\qquad(k\ge1),\]
so the terms grow without bound and the series diverges. Instead,
\[\sqrt[3]{10}=2\left(1+\dfrac14\right)^{\dfrac13}\approx2\left(1+\dfrac13\cdot\dfrac14-\dfrac19\cdot\dfrac1{16}\right)=2\cdot\dfrac{155}{144}=\dfrac{155}{72}\approx\boxed{2.1528},\]
close to the true value $2.15443\ldots$
